\section{Постановка задач}

Дана функция $f(x)$. Необходимо вычислить определённый интеграл от заданной функции на интервале $[a, b]$ с заданным параметром точности \( \varepsilon \) --- максимальной разницей между двумя итерациями \textit{(разбиениями на отрезки)}.

В \textit{варианте №2} вычисление осуществляется \textit{методом трапеций}. При реализации алгоритма необходимо учитывать следующие условия:
\begin{itemize}
  \item если у функции есть разрыв второго рода, скачок, или она не определена на какой-либо части интервала $[a, b]$, то следует установить переменные \texttt{error\_message} и \texttt{hasDiscontinuity}, а также вывести сообщение: \textit{"Integrated function has discontinuity or does not defined in current interval"};
  \item если есть устранимый разрыв первого рода, вычисление интеграла должно быть выполнено корректно;
  \item если выполнено условие $a > b$, значение интеграла должно иметь отрицательный знак.
\end{itemize}

То есть в ходе лабораторной работы будут выполнены следующие задачи:

\begin{enumerate}
  \item Разработка итерационного алгоритма численного интегрирования методом трапеций с контролем точности по заданному параметру $\varepsilon$.
  \item Реализация механизма диагностики особенностей подынтегральной функции и обработки исключительных ситуаций согласно требованиям.
  \item Организация ввода данных в формате $a$ $b$ $f$ $\varepsilon$ и формирование вывода в виде вычисленного значения интеграла $I$.
\end{enumerate}

Теоретический раздел подготовлен с опорой на внешний источник\cite{burden2011numerical}.